% ---
% §7 Conclusiones. Reemplaza las líneas 3468–3521 de Docs/Technical_Report.tex (conclusiones de TT1).
% Capítulo nuevo de TT2, escrito a partir de §6 y de la bitácora (H-031 a H-055, D-028 a D-042).
% Las conclusiones de TT1 (hallazgos del EDA y decisiones aprendidas) quedan absorbidas en §4 y §5.
% ---

\newpage
\section{Conclusiones}
\addcontentsline{toc}{section}{Conclusiones}
\label{sec:conclusiones}

Este capítulo responde a la pregunta de investigación planteada en la
sección~\ref{sec:planteamiento} a partir de los resultados del capítulo~\ref{sec:resultados}.
Primero se da la respuesta y se resume cómo se cumplió cada objetivo específico; después se explica
por qué las transformaciones cuánticas no aumentaron la separabilidad, qué aporta el trabajo y cuáles
son los límites de lo que puede afirmar. En todo el capítulo, masas y calcificaciones se tratan como
subproblemas independientes, y ``cuántico'' se refiere a simulación clásica exacta con Qiskit Aer.
